\documentclass[a4paper,10pt]{article}
\usepackage[utf8]{inputenc}
\usepackage{amssymb}
\usepackage{amsmath}
\usepackage{graphicx}
\newcommand{\HRule}{\rule{\linewidth}{0.5mm}}
\usepackage{subfigure}
\usepackage{multicol}
\usepackage[usenames,dvipsnames]{color}
\definecolor{darkgray}{rgb}{0.95,0.95,0.95}
\usepackage{listings}
\usepackage{color}
\usepackage{textcomp}
\definecolor{listinggray}{gray}{0.9}
\definecolor{lbcolor}{rgb}{0.9,0.9,0.9}
\lstset{
    backgroundcolor=\color{lbcolor},
    tabsize=4,
    rulecolor=,
    language=C++,
    basicstyle=\scriptsize,
    upquote=true,
    aboveskip={1.5\baselineskip},
    columns=fixed,
    showstringspaces=false,
    extendedchars=true,
    breaklines=true,
    prebreak = \raisebox{0ex}[0ex][0ex]{\ensuremath{\hookleftarrow}},
    frame=single,
    showtabs=false,
    showspaces=false,
    showstringspaces=false,
    identifierstyle=\ttfamily,
    keywordstyle=\color[rgb]{0,0,1},
    commentstyle=\color[rgb]{0.133,0.545,0.133},
    stringstyle=\color[rgb]{0.627,0.126,0.941},
    backgroundcolor=\color{darkgray},
}

\usepackage[left=1.0cm, right=1.0cm, top=2cm, bottom=4cm]{geometry}
\usepackage{fancyhdr}
\pagestyle{fancy}
\usepackage{lastpage}
\renewcommand\headrulewidth{1pt}
\fancyhead[L]{\textsc{Nachos Étape 3 : Multithreading}}
\fancyhead[R]{\textsc{Polytech' Grenoble}}
\renewcommand\footrulewidth{1pt}
\fancyfoot[R]{ \textsc{RICM 4}}

\begin{document}
\begin{titlepage}

\begin{center}


% Upper part of the page
\includegraphics[width=0.25\textwidth]{./images/logo}\\[1cm]

\textsc{\LARGE Polytech' Grenoble}\\[1.5cm]

\textsc{\Large RICM 4\`eme ann\'ee}\\[1.2cm]


% Title
\HRule \\[0.4cm]
{ \huge \bfseries NachOS\\[0.6cm]
Etape 3: Multithreading}
\\[0.4cm]

\HRule \\[2cm]

% Author and supervisor
\begin{minipage}{0.4\textwidth}
\begin{flushleft} \large
\emph{\'Etudiants:}\\
Elizabeth \textsc{Paz} \\
Salem \textsc{Harrache}
\end{flushleft}
\end{minipage}
\begin{minipage}{0.4\textwidth}
\begin{flushright} \large
\emph{Professeur:} \\
Vania \textsc{Marangozova}
\end{flushright}
\end{minipage}

\vfill

% Bottom of the page
{\large  février 2012}

\end{center}

\end{titlepage}

\section{Mise en place des threads \textit{utilisateurs}}

Avant de nous interesser au Multithreading, nous allons tous d'abord nous
interesser au fonctionnement des threads sous NachOS.

\subsection{Fonctionnement des threads sous NachOS}

Un thread NachOS comporte une structure de données (ContextBlock) contenant :  
\vspace{0.3cm} \\
\begin{tabular}{|l|l|}
\hline 
 \textbf{status} & JUST\_CREATED, RUNNING, READY, BLOCKED \\ \hline 
 \textbf{stack}& L'emplacement (bas) de la pile \\ \hline
 \textbf{stackTops}& Sommet de la pile \\ \hline
 \textbf{machineState}&Le contenu de l'ensemble de registres \\ \hline
 \textbf{userRegisters} &Le contenu de chaque registre (Thread user) \\ \hline
 \textbf{space}& L'espace d'adressage (Thread user) \\ \hline
\end{tabular} \\
\vspace{0.2cm} \\

Un thread est tout d'abord créé avec la méthode Thread() dans le fichier 
\textit{threads/thread.cc}. Cette méthode initialise la pile, le pointeur de la 
pile et le thread est en état: \textit{JUST\_CREATED}. Cet état temporaire indique
que le thread est créé mais n'est pas encore prêt (READY) à être executé. 

L'appel Fork (différent du Fork Posix) qui va rendre ce
thread exécutable. Concrètement, la pile de notre thread sera alloué et 
initialisé dans cette fonction. Il va être ensuite mise dans l'ordannanceur de NachOS
pour qu'il puisse être executé. La pile d'un thread NachOS

La fonction \textit{saveUserState} : sert à sauvegarder le contexte d'exécution
du programme lors d'une commutation. Les informations telles que l'état des
registres ou encore le pointeur de pile sont sauvegardées dans le contextBlock.
Ce contexte sera restauré lorsque le Thread est élu avec la méthode
\textit{restoreUserState}.

\subsection{Description de l'instalation d'un programme utilisateur dans
l'environnement NachOS}

Le lancement pas à pas permet de mettre en evidence plusieurs étapes lors du
lancement d'un programme dans l'environnement Nachos.

\begin{lstlisting}
Initializing address space, num pages 12, size 1536
Initializing code segment, at 0x0, size 384
Initializing data segment, at 0x180, size 32
Initializing stack register to 1520
Starting thread "main" at time 10
\end{lstlisting}

\subsubsection{Création de l'espace d'adressage}

La premiere étape est la construction de l'espace d'adressage dans
\textit{::AddrSpace}. Le fichier exécutable est stocké dans un objet Noff. Ce
dernier contient le text, les données et les variables non initialisées. La
taille total de l'espace d'adressage est la somme de ces trois zones au quelle
on ajoute la taille de pile utilisateur qui est de 4Ko (1024 * 4).

\begin{lstlisting}
#define UserStackSize		1024
\end{lstlisting}

Avant de stocker cet espace d'adressage il faut connaitre le nombre de pages
necessaires. Le système calcule le nombre de pages puis les initialise pour
avoir la correspondance entres adresses virtelles et adresses physiques. Une
fois les pages correctement initialisées, on copie les differentes zones en
mémoires et ferme le fichier exécutable.

\subsubsection{Initialisation des registres}

Dans \textit{::InitRegisters} 
\begin{lstlisting}
for (i = 0; i < NumTotalRegs; i++)
   machine->WriteRegister (i, 0);
machine->WriteRegister (PCReg, 0);
machine->WriteRegister (NextPCReg, 4);
machine->WriteRegister (StackReg, numPages * PageSize - 16)
\end{lstlisting}

Comme on peut le voir, tous les registres sont mis à zéro. Trois
registres sont ensuite initialisés : PCReg qui pointe sur la première
instruction du programme (0x0). MIPS utilise également un pointeur sur
l'instruction N+1, soit à l'adresse 0x4. Enfin le pointeur de pile (StackReg)
qui pointe vers la fin de l'espace d'adressage.

[...Pourquoi -16 ?...]

\subsubsection{Lancement du programme utilisateur}

L'exécution du programme est effectuée par la méthode \textit{::Run}. Celle ci
parcours les inscrutions en incrementant l'horloge à l'aide de la fonction
OneTick.

\subsection{Threads utilisateurs: étape préliminaire}

ù\subsubsection{Appels systèmes}

  La fonction \textit{UserThreadCreate(...)}

\subsubsection{Création de la fonction}

Sans l'option -rs fait le thread  main ne donne JAMAIS la main au thread user!


\section{Plusieurs threads par processus}


\end{document}
